\section{Yöntem}
\label{sec:yontem}

Bu bölüm beş alt bölümden oluşmaktadır: deney ortamı ve görev tanımı, tüm
koşuların paylaştığı simülasyon kurulumu, arazi temsili ve karıştırma
probleminin matematiksel formülasyonu, analitik (CPU) değerlendirme
protokolü ve gerçek Isaac Lab/Isaac Sim eğitim protokolü. Üç karıştırma
yöntemi de önce analitik/CPU düzeyinde, ardından gerçek GPU paralel fizik
simülasyonu içinde aynı sabit geçiş yarıçapı ve aynı hard-truncation
kuralı altında değerlendirilmiştir; bu tutarlılık, iki değerlendirme
katmanı arasındaki bulguların karşılaştırılabilir olmasını sağlamaktadır.

\begin{figure}[t]
\centering
\includegraphics[width=\linewidth]{fig01_methodology.png}
\caption{Genel metodoloji akış şeması: arazi üretiminden (DEM + prosedürel
ince detay), karıştırma fonksiyonu seçimine, analitik/CPU
değerlendirmesine ve gerçek Isaac Lab/SAC eğitim-analiz hattına.}
\label{fig:methodology}
\end{figure}

\subsection{Deney Ortamı ve Görev Tanımı}
\label{sec:ortam}

Deney ortamı, Isaac Lab doğrudan iş akışı (direct-workflow) API'si
üzerinde tanımlanmış bir görev olan \texttt{LunarRocket-Lander-Direct-v0}
'dır. Ajan, dört bacaklı ve iki eksenli gimbal itkili bir roketi temsil
eder; gözlem uzayı konum, hız, yönelim, açısal hız, IMU ve altimetre
okumaları ile hedefe göreli konumu içerir; eylem uzayı ana itki
büyüklüğünü ve gimbal açısını kapsar. Ajan, Stable-Baselines3 (SB3) SAC
(Soft Actor-Critic) algoritmasıyla, NVIDIA PhysX üzerinde GPU paralel
olarak çalışan yüzlerce ortam örneğinde eş zamanlı eğitilmektedir.

Arazi, tek bir prosedürel adımda değil, iki katmanlı bir temsille üretilir:
(i) gerçek bir NASA dijital yükseklik modelinden (DEM) türetilen, rastgele
eğim, krater ve kaya alanlarıyla zenginleştirilmiş bir küresel (global)
makro yüzey; (ii) ajanın başlangıç ve hedef konumlarının yakınında,
santimetre ölçekli çok frekanslı bir gürültü fonksiyonuyla üretilen
yüksek çözünürlüklü bir yerel (local) regolit deseni. Bu iki katman, ortam
sıfırlaması (reset) sırasında arka planda çalışan bir iş parçacığı havuzu
tarafından önceden hazırlanmış bir ``terrain havuzu'' içinden bir yuvanın
(slot) kopyalanmasıyla GPU'ya yüklenir; aktif bölümler (episode) hiçbir
zaman bölüm ortasında değiştirilmez, yeni üretimler yalnızca sıfırlama
anlarında devreye girer. Eğitim, spawn/hedef mesafesi ve komutlanabilir
gimbal yetkisinin, kayan pencereli başarı oranına göre kademeli arttığı
bir müfredat (curriculum) düzeni içinde yürütülmektedir.

Bu iki katmanın birleştiği yerde, yerel yamanın ince detayı, küresel
yüzeye sabit bir yarıçap içinde ``söndürülerek'' (fade) eklenir; bu
söndürme fonksiyonunun matematiksel şekli, bu çalışmanın konusudur ve
\cref{sec:blend}'te tanımlanmaktadır.

\Cref{fig:scene-wide} ve \cref{fig:scene-closeup}, eğitilmiş bir politikanın
gerçek Isaac Lab/Isaac Sim ortamında çalıştırılmasından alınmış, doğrudan
simülasyondan yakalanmış görüntülerdir (offscreen render,
\texttt{ISAACLAB\_CHECKPOINT} ile yüklenen bir kontrol noktası üzerinden;
araç iniş yaklaşımının son evresinde, zemine yaklaşık 0,6--0,7~m
yükseklikteyken yakalanmıştır). Bu yakalama için iki değişiklik
yapılmıştır: (i) sahnenin varsayılan gölgesiz dome light'ı kısılmış ve
düşük açılı, gerçek gölge düşüren yönlü bir ``güneş'' ışığı eklenmiştir,
böylece arazi kabartması ve gölgeler görünür hâle gelmiştir; (ii) iniş
bölgesi, 128x128 küresel yüzeyin yanı sıra 512x512 çözünürlüklü gerçek
yüksek-detay yerel yama (bkz.\ \cref{sec:ortam}, ``yerel regolit deseni'')
ile render edilmiştir, böylece santimetre ölçekli ince doku (toz/kum
taneciği düzeyinde pürüzlülük) da görünür olmuştur. Aracın gerçekten dik
durduğu varsayıma değil, simülasyonun kendi gövde-yönelim verisine
dayanmaktadır: bu koşu boyunca gövdenin yerel yukarı ekseninin dünya $+z$
ekseniyle hizası ($\mathrm{up}_z$) sabit biçimde $+1{,}000$ ölçülmüştür,
yani araç devrilmiş değil, tam dikeydir. \Cref{fig:scene-wide}, dört
bacaklı iniş aracını gerçek kraterli/kayalık ay zemini üzerinde
göstermektedir; sırt hattı, taneli/pürüzlü yüzey dokusu, gölgeler ve
dağınık kayalar, \cref{sec:ortam}'da tanımlanan DEM tabanlı makro
yüzeyin ve prosedürel kaya/krater/ince-detay alanlarının çalışır
durumdaki hâlidir. \Cref{fig:scene-closeup}, aynı karedeki iniş aracının
yakın çekimidir; gövde, iki eksenli gimbal itki ekseninin uçları ve
zeminin taneli dokusu birlikte görülebilmektedir. Bu görüntüler yalnızca
ortamı görsel olarak belgelemek amacıyla sunulmuştur; belirli bir iniş
sonucunu (başarılı/başarısız) temsil ettiği iddia edilmemektedir.

\begin{figure}[t]
\centering
\includegraphics[width=\linewidth]{fig02_scene_wide.png}
\caption{Eğitilmiş politikanın gerçek Isaac Sim ortamında
çalıştırılmasından alınmış sahne görüntüsü: dört bacaklı iniş aracı,
kraterli/kayalık ve taneli dokulu ay zemini üzerinde (yüksek-detay yerel
yama + belgeleme amaçlı yönlü ışık; araç simülasyon verisiyle doğrulanmış
biçimde tam dikeydir, $\mathrm{up}_z=+1{,}000$).}
\label{fig:scene-wide}
\end{figure}

\begin{figure}[t]
\centering
\includegraphics[width=0.6\linewidth]{fig03_scene_closeup.png}
\caption{Aynı sahnenin iniş aracına yakınlaştırılmış görünümü: gövde,
gimbal itki ekseninin uçları ve zeminin santimetre ölçekli taneli dokusu
birlikte görülebilmektedir.}
\label{fig:scene-closeup}
\end{figure}

\subsection{Simülasyon Kurulumu}
\label{sec:simkurulum}

Bu alt bölüm, hem analitik (CPU) protokolün hem de gerçek Isaac
Lab/Isaac Sim eğitim koşularının paylaştığı simülasyon yapılandırmasını
tek bir yerde toplamaktadır (\cref{tab:config}). Fizik, NVIDIA PhysX
üzerinde GPU-paralel olarak $1/120$ saniyelik sabit adımla
(\SI{120}{\hertz}) çözülmekte; politika, \texttt{decimation}~$=2$ ile
\SI{60}{\hertz}'de eylem üretmekte ve her bölüm \SI{20}{\second}, yani
1200 kontrol adımı sürmektedir. Ay yerçekimi
(\SI{1.62}{\meter\per\second\squared}), düşük restitüsyonlu regolit
benzeri bir temas malzemesi ve 512 eş zamanlı ortam örneği tüm koşularda
sabit tutulmuştur.

Araç ve tahrik parametreleri (kütle, itki zarfı, fiziksel ve politika
gimbal limitleri, TVC aktüatör kazançları) ile SAC hiperparametreleri, üç
karıştırma yöntemi arasında değiştirilmemiştir. Yöntemden yönteme değişen
tek büyüklük, \cref{sec:blend}'te tanımlanan ince-detay söndürme
fonksiyonu ve ona bağlı parametrelerdir (\cref{tab:config}'in son grubu;
\texttt{ISAACLAB\_TERRAIN\_BLEND\_METHOD} ortam değişkeniyle seçilir). Bu
sabitleme, \cref{sec:bulgular}'da gözlemlenen farkların yalnızca
karıştırma yöntemine atfedilebilmesi için gereklidir.

Gerçek eğitim karşılaştırması, hesaplama bütçesi kısıtı nedeniyle tam
ölçekli (1200 iterasyon) yapılandırmanın kısaltılmış bir versiyonuyla
-yöntem başına üç bağımsız tohum, koşu başına yaklaşık 5,73 milyon ortam
adımı- yürütülmüştür; bu ölçek sınırlaması \cref{sec:egitimprotokolu} ve
\cref{sec:egitimsonuclari}'te ayrıca ele alınmaktadır.

\begin{table}[t]
\centering
\caption{Analitik protokol ile gerçek Isaac Lab/Isaac Sim eğitim
koşularının paylaştığı simülasyon, araç, RL eğitim ve arazi karıştırma
parametreleri. Yalnızca son grup (karıştırma parametreleri), yöntemden
yönteme değişir.}
\label{tab:config}
\footnotesize
\begin{tabular}{@{}p{0.19\linewidth}p{0.33\linewidth}p{0.38\linewidth}@{}}
\toprule
\textbf{Grup} & \textbf{Parametre} & \textbf{Değer} \\
\midrule
\multirow{10}{*}{Fizik / simülasyon}
 & Fizik motoru & NVIDIA PhysX 5 (GPU ardışık düzeni) \\
 & Simülasyon adımı (dt) & $1/120$~s (\SI{120}{\hertz}) \\
 & Kontrol frekansı (decimation = 2) & \SI{60}{\hertz} \\
 & Bölüm süresi & \SI{20}{\second} (1200 kontrol adımı) \\
 & Yerçekimi & $0,0,-1.62$~m/s$^2$ \\
 & Paralel ortam sayısı & 512 \\
 & Ortam aralığı (env\_spacing) & \SI{90}{\meter} \\
 & Fizik kopyalama & replicate\_physics + clone\_in\_fabric (açık) \\
 & Sürtünme (statik / dinamik) & 0.85 / 0.65 (combine: multiply) \\
 & Restitüsyon & 0.02 \\
\midrule
\multirow{9}{*}{Araç ve tahrik}
 & Araç kütlesi & 1.66~kg (1.63 gövde + gimbal bağlantıları) \\
 & Maksimum itki & 30~N \\
 & Maksimum komut throttle & 0.35 ($\approx$~10.5~N etkin) \\
 & Fiziksel gimbal limiti & $\pm25^\circ$ (mekanik sert durdurucu) \\
 & Politika gimbal zarfı & $4^\circ$--$8^\circ$ (müfredat zorluğuyla ölçeklenir) \\
 & İtki kaldıraç kolu & 0.25~m \\
 & TVC aktüatör (implicit) & effort 2.5, stiffness 10, damping 0.1, armature 0.002 \\
 & Gözlem / eylem boyutu & 125 / 3 \\
 & Başlangıç irtifası & 25~m (dik yönelim) \\
\midrule
\multirow{12}{*}{RL eğitimi (SB3 SAC)}
 & Politika ağı & MLP [384, 256, 128], ELU \\
 & Replay buffer & 1\,000\,000 \\
 & learning\_starts & 16\,384 \\
 & Batch boyutu & 1024 \\
 & Öğrenme oranı & $3\times10^{-4}$ \\
 & $\gamma$ / $\tau$ & 0.99 / 0.005 \\
 & train\_freq / gradient\_steps & 1 / 32 \\
 & Entropi katsayısı & otomatik (target\_entropy = $-1.5$) \\
 & Normalizasyon & normalize\_input + normalize\_value açık, clip\_obs = 10 \\
 & Karşılaştırma bütçesi & 350 iter $\times$ 32 $\times$ 512 $\approx$ 5.73~M ortam adımı \\
 & Tohumlar & 42, 7, 123 (yöntem başına $n=3$, ardışık) \\
 & Cihaz & tek GPU (cuda:0) \\
\midrule
\multirow{9}{*}{Karıştırma parametreleri}
 & Geçiş yarıçapı ($R$) & 8.0~m \\
 & Yöntem seçenekleri & smoothstep / gaussian / hybrid (ISAACLAB\_TERRAIN\_BLEND\_METHOD) \\
 & Kalibre Gauss katsayısı ($k$) & $\ln(1000)\approx 6.9078$ \\
 & Hibrit temas yarıçapı ($r_c$) & 1.5~m \\
 & Hibrit switch genişliği ($w$) & 0.3~m \\
 & İnce-detay genlik aralığı & 0.006--0.020~m \\
 & Makro yükseklik / eğim ölçeği & 0.45~m / 0.035 \\
 & Krater / kaya sayısı (bölüm başına) & 3 / 5 \\
 & Terrain havuzu (boyut / yenileme) & 15 yuva / 60 atama \\
\bottomrule
\end{tabular}
\end{table}

\subsection{Arazi Temsili ve Karıştırma Fonksiyonları}
\label{sec:blend}

Yerel yamanın ince detay yüksekliği, küresel yüzeye üç aday söndürme
(fade) fonksiyonundan biriyle eklenir. Her üç fonksiyon da $r$ merkezden
uzaklık ve $R=\texttt{terrain\_detail\_radius\_m}=8{,}0$~m sabit geçiş
yarıçapı olmak üzere $r\ge R$ için tam olarak sıfıra zorlanır (hard
truncation); bu, üç yöntemi aynı geçiş bütçesi altında karşılaştırılabilir
kılan temel metodolojik düzenlemedir ve üretim kodunda da bu şekilde
uygulanmaktadır.

\textbf{Smoothstep} (kompakt destek, varsayılan/mevcut yöntem): $t =
\mathrm{clamp}(1-r/R,\,0,\,1)$ olmak üzere
\begin{equation}
\mathrm{fade}(r) = t^2\,(3-2t).
\label{eq:smoothstep}
\end{equation}
Bu fonksiyon $r=R$'de zaten tam olarak sıfırdır; ek bir kesme adımına
ihtiyaç duymaz ve C1 sürekliliği analitik olarak garantidir.

\textbf{Gauss} (sonsuz destek, önceden kullanılan yöntem):
\begin{equation}
\mathrm{fade}(r) = \exp\!\left(-k\,(r/R)^2\right).
\label{eq:gauss}
\end{equation}
Bu çalışmada iki $k$ değeri ayrı ayrı değerlendirilmiştir: önceden
kullanılan (as-shipped) sabit $k=1{,}0$ ve bu çalışmada önerilen, Gauss'un
pratik desteğinin $R$'de yaklaşık binde bir düzeyine inmesini sağlayan
kalibre edilmiş sabit $k=\ln(1000)\approx 6{,}9078$. Üretim kodundaki
\texttt{cfg.terrain\_gaussian\_k} alanı, varsayılan olarak kalibre
edilmiş sabiti kullanır; üretim öncesi $k=1{,}0$ yalnızca bu çalışmanın
analitik/CPU karşılaştırmasında, mevcut kusuru göstermek amacıyla ayrı
bir varyant (\texttt{gaussian\_shipped}) olarak test edilmiştir. Her iki
$k$ değeri için de fonksiyon, $r\ge R$'de dışarıdan sıfıra zorlanır.

\textbf{Hibrit} (önerilen yöntem): İniş ayağı temas izi yarıçapı $r_c =
\texttt{terrain\_hybrid\_contact\_radius\_m}=1{,}5$~m ve switch genişliği
$w=\texttt{terrain\_hybrid\_switch\_width\_m}=0{,}3$~m olmak üzere,
$[r_c-w/2,\,r_c+w/2]$ aralığında smoothstep ağırlıklı bir $\mathrm{switch}(r)$
tanımlanır:
\begin{equation}
u = \mathrm{clamp}\!\left(\frac{\mathrm{hi}-r}{\mathrm{hi}-\mathrm{lo}},\,0,\,1\right),
\qquad
\mathrm{switch}(r) = u^2\,(3-2u),
\label{eq:switch}
\end{equation}
burada $\mathrm{lo}=r_c-w/2$, $\mathrm{hi}=r_c+w/2$. Nihai söndürme,
\begin{equation}
\mathrm{fade}(r) = \mathrm{switch}(r)\cdot\mathrm{gauss}(r) +
\left(1-\mathrm{switch}(r)\right)\cdot\mathrm{smoothstep}(r)
\label{eq:hybrid}
\end{equation}
olarak tanımlanır; yani ayak temas bölgesine yakın kısımda ağırlıklı
olarak Gauss, ötesinde ağırlıklı olarak smoothstep kullanılır ve iki
fonksiyon arasındaki geçiş, kendisi de C1 düzgün olan bir smoothstep
ağırlıkla yumuşatılır. $r=R$ sınırında $\mathrm{switch}(r)=0$ olduğundan,
hibrit yöntem bu sınırda analitik olarak smoothstep ile birebir aynıdır.

Bu üç yöntem, üretim kodunda \texttt{ISAACLAB\_TERRAIN\_BLEND\_METHOD}
ortam değişkeniyle (\texttt{smoothstep} \textbar\ \texttt{gaussian}
\textbar\ \texttt{hybrid}) seçilebilir şekilde, mevcut varsayılan
davranışı (smoothstep) değiştirmeden entegre edilmiştir.

\subsection{Analitik (CPU) Değerlendirme Protokolü}
\label{sec:analitikprotokol}

Analitik değerlendirme, Isaac Lab/torch bağımlılığı olmadan, üretim
kodundaki formüllerin bire bir NumPy taşıması üzerinde, iki bağımsız
izleme (track) ile yürütülmüştür: (i) mesh-track, gerçek üçgen ağı üretip
vertex sayısı ve süre ölçen bir izleme; (ii) analytic-track, yükseklik
alanını doğrudan sürekli bir fonksiyon olarak sorgulayan bir izleme. İki
izlemenin birbirini doğrulaması, bulguların mesh ayrıklaştırmasına özgü
bir yapay etki olmadığını göstermek için kullanılmıştır.

Süreklilik metrikleri, sonlu farklarla hesaplanmıştır: C0 (yükseklik)
sürekliliği, geçiş sınırının hemen iki tarafındaki noktalar arasındaki
yükseklik farkının mutlak değeri (uzamsal epsilon = 0,05~m); C1 (eğim)
sürekliliği, aynı noktalardaki yerel gradyanın açısal farkı (gradyan
epsilon = 0,02~m) olarak tanımlanmıştır. Ölçümler 720 açısal örnek
üzerinden alınmış ve hem maksimum hem RMS (kök ortalama kare) değerleri
raporlanmıştır. İlk ölçümlerde, birleşik (küresel + yerel) yükseklik
üzerinden yapılan ölçümün, doğal arazi eğiminin epsilon aralığındaki
varyasyonuyla karışarak asıl dikiş sinyalini bastırdığı görülmüş; bu
nedenle yalnızca söndürülmüş ince detay bileşenini
($\mathrm{fade}(r)\cdot\mathrm{ince\_detay}(x,y)$) izole eden bir
``detail-only'' ölçüm modu tanımlanmış ve bu, çalışmanın başlık
(headline) metriği olarak benimsenmiştir; birleşik yüzey üzerindeki ölçüm
ikincil/gerçekçi bir referans olarak korunmuştur.

Bellek maliyeti, vertex başına 32 bayt ve çarpışma (collision) ağı için
2x çarpan varsayımıyla, iki-mesh (literal, ayrı küresel ve yerel ağlar)
ile tek-grid (baked, tek birleşik ağ) muhasebe biçimleri için ayrı ayrı
tahmin edilmiştir. Sorgu verimi (throughput), 512 ortam ve ortam başına
64 küresel + 24 yerel + 4 ince-detay sorgusu (toplam 92 sorgu/ortam)
varsayımıyla, saf söndürme formülünün (fade-only, terrain örneklemesinden
izole) mikro-saniye cinsinden maliyeti ayrıca ölçülmüştür.

Ana karşılaştırma, dört yöntem varyantı (smoothstep, üretim öncesi Gauss,
kalibre edilmiş Gauss, hibrit) için otuz bağımsız rastgele tohum
üzerinde tekrarlanmıştır; istatistiksel anlamlılık, her tohumda
hesaplanan $\log_{10}$ oran farkları üzerinde Wilcoxon işaretli sıra
testiyle değerlendirilmiştir (parametrik testler yerine bu test tercih
edilmiştir, çünkü bazı karşılaştırmalarda -örneğin hibrit ile
smoothstep, $r=R$'de analitik olarak özdeş olduğundan- fark dağılımı
sıfır varyansa yakınsamakta ve Cohen's $d$ gibi parametrik büyüklük
ölçütleri yozlaşmaktadır).

Bulguların sağlamlığı dört ayrı ablasyon analiziyle sınanmıştır: (i)
temas yarıçapı ($r_c$) taraması, 1,0--4,0~m aralığında; (ii) switch
genişliği ($w$) taraması, 0,05--2,0~m aralığında; (iii) geçiş yarıçapı
($R$) taraması, 4--16~m aralığında, kusurun ölçek değişmezliğini test
etmek için; (iv) grid çözünürlüğü taraması, yaklaşık 0,31--1,27~m grid
aralığı üzerinde, log-log doğrusal regresyonla ampirik yakınsama derecesi
(convergence order) tahmin edilerek.

\subsection{Gerçek Isaac Lab/Isaac Sim Eğitim Protokolü}
\label{sec:egitimprotokolu}

Analitik değerlendirmenin ardından, üç yöntem (smoothstep, kalibre
edilmiş Gauss, hibrit) gerçek bir Isaac Lab/Isaac Sim eğitim ortamında
karşılaştırılmıştır. Bu, çalışmanın özgün değerinin merkezini
oluşturur: karşılaştırma soyut bir sayısal deney olarak değil, gerçek
kraterli/eğimli ay zemini üzerinde, GPU paralel PhysX fizik simülasyonu
içinde çalışan bir SAC eğitim döngüsüyle gerçekleştirilmiştir.

Her yöntem, \texttt{ISAACLAB\_TERRAIN\_BLEND\_METHOD} ortam değişkeniyle
seçilerek, aynı donanımda (tek makine, tek GPU), aynı hiperparametre
yapılandırmasıyla (\texttt{sb3\_sac\_cfg.yaml}) ve GPU/VRAM çakışmasını
önlemek amacıyla eş zamanlı değil ardışık olarak eğitilmiştir. Her koşu
5.734.400 ortam adımı (yaklaşık 350 iterasyon) sürmüştür; bu, hesaplama
bütçesi kısıtı nedeniyle tam ölçekli (1200 iterasyon) yapılandırmanın
kısaltılmış ama anlamlı bir sinyal verecek şekilde küçültülmüş bir
versiyonudur ve bu, çalışmanın açıkça belirtilmesi gereken bir ölçek
sınırlamasıdır.

İlk ölçüm, SAC ajanının yapılandırma dosyasındaki varsayılan tohumla
(seed~$=42$) yöntem başına tek koşu olarak yapılmıştır. Bu tek-koşu
tasarımının Henderson ve ark.'nın \citep{henderson2018deep} uyardığı
türden bir tuzak olduğu göz önünde bulundurularak, çalışma iki ek
bağımsız tohumla (seed~$=7$ ve seed~$=123$) genişletilmiş ve her yöntem
için toplam üç bağımsız koşu ($n=3$) elde edilmiştir
-\texttt{ISAACLAB\_SEED} ortam değişkeni bu amaçla üretim koduna
eklenmiştir. GPU paralel PhysX fizik simülasyonunun kendi içsel
belirlenimsizliği (iş parçacığı zamanlaması, kayan nokta indirgeme
sırası) ve terrain havuzunun arka planda çalışan CPU üretim sürecinin
zamanlamaya bağlı olması nedeniyle, aynı tohumla bile iki koşu birbirinin
bit eşdeğeri tekrarı değildir; dolayısıyla üç tohum, ajan başlatma
kaynaklı varyansın yanı sıra bu sistem-düzeyi belirlenimsizliği de
örnekler.

Eğitim metrikleri (başarı oranı, zaman aşımı oranı, sert iniş oranı,
müfredat zorluğu, entropi katsayısı, xy-ilerleme ödülü) TensorBoard
üzerinden loglanmış ve olay dosyaları (event files)
\texttt{tensorboard.backend.event\_processing.EventAccumulator} ile
ayrıştırılmıştır. Her (yöntem, tohum) koşusu için nihai performans, o
koşunun son yüzde onluk penceresindeki loglanmış noktaların ortalaması
alınarak TEK bir özet değere indirgenmiştir; yöntemler arası
karşılaştırmalar bu üç bağımsız özet değer üzerinde (bağımsız örneklem
olarak) Mann-Whitney U testi ve 5000 tekrarlı bootstrap \%95 güven
aralığı ile yapılmıştır. Bu, bir tek koşunun içindeki otokorelasyonlu
(bağımsız olmayan) noktaları sahte-bağımsız örneklem gibi kullanmaktan
(pseudoreplication) kaçınan, istatistiksel olarak doğru birimdir; ancak
$n=3$ ile istatistiksel güç düşüktür ve bu husus \cref{sec:egitimsonuclari}
ve \cref{sec:tartisma}'te açıkça vurgulanmaktadır.
